\capitulo{5}{Aspectos relevantes del desarrollo del proyecto}

\section{Ciclo de vida del proyecto}

El proyecto se ha organizado en cinco sprints siguiendo una metodología ágil.

\begin{itemize}
    \item Sprint 1: Análisis del problema empresarial, definición de requisitos y arquitectura inicial.
    \item Sprint 2: Provisionamiento de infraestructura Azure y configuración de Terraform.
    \item Sprint 3: Desarrollo de la API Flask, portal web y persistencia en Azure Storage.
    \item Sprint 4: Seguridad con Key Vault, Managed Identity, portal web y despliegue en Azure.
    \item Sprint 5: Pruebas, calidad SonarCloud, documentación y validación del prototipo.
\end{itemize}

\section{Arquitectura del sistema}

La arquitectura implementada sigue un modelo de servicios gestionados en Microsoft Azure. Los componentes principales son:

\begin{itemize}
    \item \textbf{Azure App Service}: hospeda la API Flask y el portal web empresarial en Python 3.11.
    \item \textbf{Azure Storage}: persiste las solicitudes operativas TI en un contenedor Blob privado.
    \item \textbf{Azure Key Vault}: almacena el token de autenticación de la API.
    \item \textbf{Azure Logic Apps}: automatiza la entrada de solicitudes externas mediante un trigger HTTP.
    \item \textbf{Application Insights y Azure Monitor}: permiten revisar telemetría, métricas y comportamiento operativo del servicio desplegado.
    \item \textbf{Scripts de despliegue}: automatizan la configuración y publicación en Azure mediante Azure CLI.
\end{itemize}

El flujo principal del sistema es el siguiente:

\begin{enumerate}
    \item Un empleado envía una solicitud operativa TI desde el portal web, mediante la API REST o a través de la Logic App.
    \item Azure App Service procesa la petición con Gunicorn.
    \item La API valida los datos y clasifica automáticamente la solicitud.
    \item El sistema genera una recomendación operativa para el equipo de TI.
    \item Los datos se almacenan en Azure Blob Storage.
    \item Los secretos se obtienen desde Azure Key Vault mediante Managed Identity.
    \item El endpoint de métricas y Application Insights permiten observar el estado agregado del sistema.
\end{enumerate}

\begin{figure}[h]
\centering
\includegraphics[width=0.9\textwidth]{arquitectura_final_azure.png}
\caption{Arquitectura final del prototipo desplegado en Microsoft Azure. Fuente: elaboración propia.}
\label{fig:arquitectura-general-memoria}
\end{figure}

La figura \ref{fig:infografia-arquitectura-azure} resume visualmente los servicios cloud utilizados y su relación con el proceso empresarial automatizado.

\begin{figure}[h]
\centering
\includegraphics[width=0.9\textwidth]{infografia_de_arquitectura_Azure_v2.png}
\caption{Infografía de arquitectura cloud del prototipo. Fuente: elaboración propia.}
\label{fig:infografia-arquitectura-azure}
\end{figure}

\section{API REST implementada}

La API desarrollada en \texttt{src/app.py} expone los siguientes endpoints:

\begin{itemize}
    \item \texttt{GET /}: estado del servicio y versión.
    \item \texttt{GET /portal}: portal web para solicitudes operativas TI.
    \item \texttt{GET /health}: comprobación de salud y modo de almacenamiento.
    \item \texttt{POST /solicitudes}: creación de solicitudes con clasificación automática.
    \item \texttt{GET /solicitudes}: listado con filtros por estado, prioridad y tipo.
    \item \texttt{GET /metricas}: métricas agregadas del sistema.
\end{itemize}

La tabla \ref{tabla:endpoints_api} resume el contrato principal de la API, indicando el mecanismo de seguridad y la forma en que se ha validado cada operación.

\begin{table}[h]
\centering
\small
\begin{tabularx}{\textwidth}{l l X X}
\textbf{Método} & \textbf{Ruta} & \textbf{Seguridad} & \textbf{Validación realizada} \\
\hline
GET & \texttt{/} & Pública & Comprobación de estado general del servicio. \\
GET & \texttt{/portal} & Pública & Acceso al portal web desplegado en App Service. \\
GET & \texttt{/health} & Pública & Verificación de estado y modo de almacenamiento. \\
POST & \texttt{/solicitudes} & Pública con validación de entrada & Creación de solicitud y clasificación automática. \\
GET & \texttt{/solicitudes} & Token Bearer & Consulta protegida con filtros y paginación. \\
GET & \texttt{/metricas} & Token Bearer & Revisión de métricas agregadas del sistema. \\
\end{tabularx}
\caption{Endpoints principales de la API}
\label{tabla:endpoints_api}
\end{table}

Al crear una solicitud, el módulo \texttt{classifier.py} analiza el título y la descripción para asignar tipo de solicitud (\textit{acceso}, \textit{entorno}, \textit{aplicación}, \textit{configuración}, \textit{incidencia}), prioridad (\textit{baja}, \textit{media}, \textit{alta}), categoría (\textit{seguridad}, \textit{infraestructura}, \textit{soporte}, \textit{general}) y una recomendación operativa, devolviendo un valor de confianza asociado.

\section{Componente de inteligencia artificial}

Se ha incorporado un clasificador ligero basado en análisis semántico de palabras clave y reglas de negocio. Este componente cumple el objetivo del TFG de aportar análisis automático de información operativa sin depender de infraestructuras corporativas complejas. El clasificador evalúa términos relacionados con seguridad, infraestructura y soporte para determinar la categoría, prioridad y recomendación de cada solicitud operativa TI.

\section{Automatización del despliegue}

Se ha implementado un proceso de despliegue reproducible mediante \texttt{scripts/deploy-azure.ps1}. Este script:

\begin{enumerate}
    \item Comprueba la sesión de Azure CLI.
    \item Obtiene la configuración de Azure Storage y Key Vault.
    \item Habilita Managed Identity en Azure App Service.
    \item Concede a la aplicación permisos sobre el secreto almacenado en Key Vault.
    \item Configura variables de entorno y publica el código Python en Azure App Service.
\end{enumerate}


\section{Automatización empresarial con Logic App}

La Logic App actúa como capa de orquestación para integrar la plataforma con otros sistemas de la empresa. En el prototipo, un sistema externo puede enviar una petición HTTP a la Logic App y esta invoca de forma controlada el endpoint \texttt{POST /solicitudes} de la API desplegada en App Service. La API conserva la responsabilidad de validar datos, clasificar la solicitud, generar la recomendación y persistir el registro en Azure Blob Storage.

Este flujo resulta útil para un escenario empresarial realista: un formulario corporativo, una herramienta de soporte, un buzón automatizado o un sistema de atención interna podrían crear solicitudes sin acceder manualmente al portal. La Logic App no sustituye al equipo de TI ni al backend, sino que automatiza la entrada de trabajo y permite que el equipo técnico consulte, priorice y supervise las solicitudes desde la plataforma.
\section{Seguridad y gestión de secretos}

La autenticación de la API utiliza un token Bearer almacenado en Azure Key Vault. La App Service accede al secreto mediante Managed Identity, evitando credenciales embebidas en el código fuente. El almacenamiento de solicitudes reside en un contenedor Blob con acceso privado.

\section{Validación del prototipo}

Las pruebas unitarias en \texttt{tests/test\_api.py} verifican la creación de solicitudes, el listado, las métricas, el portal web y la clasificación automática. El prototipo puede ejecutarse localmente con almacenamiento en fichero JSON o desplegarse en Azure con persistencia en Blob Storage.

